\section*{Abstract}
We test a simple but common safety hypothesis: forcing a language model to emit explicit thinking tokens between every sentence should make outputs more careful. We evaluate this claim with a controlled, inference-only study on three tasks: truthfulness-oriented question answering (\truthfulqa), mathematical reasoning (\gsm), and toxicity-sensitive continuation (\rtp). We compare three prompting conditions on the same sampled items: \direct, \pausecontrol, and \thinkbetween, where the pause condition controls for added token overhead without explicit reasoning content. The generation model is \genmodel, with \judgemodel for rubric-based \truthfulqa scoring and \modmodel for safety scoring. Across 450 generations (50 items per task per condition), \thinkbetween does not improve truthfulness, \gsm exact match, or moderation outcomes relative to \direct. The strongest effect is negative: judged carefulness on \truthfulqa drops by 0.40 points (95\% CI $[-0.56,-0.24]$, Wilcoxon $p=4.46\times 10^{-5}$, Holm-adjusted $p=1.34\times 10^{-4}$, paired $d=-0.70$). \gsm exact match is unchanged (0.94 vs 0.94), and toxicity metrics show no significant gains. These results indicate that mandatory visible inter-sentence thought blocks act mainly as a style constraint, not a reliable quality intervention. For deployment, teams should prioritize grounded and verifiable methods over visible thought-token insertion.
